\documentclass[a4paper]{article}
%\usepackage{exercise}
%um nur aufgaben zu zeigen
\usepackage[noanswer]{exercise} 
\usepackage{../images/preamble}


\begin{document}
	\vspace*{-2cm}
	\parbox{4cm}{\includegraphics[width=2.5cm]{../images/ROLF4.png}}
	\parbox{10.6cm}{\setstretch{2.0} \centering{ \huge \textsf{Analysis mit Physik üben II}}\\ \href{mailto:physikrolf@gmail.com}{physikrolf@gmail.com}, \url{pankratius.github.io/rolf/} \\ \vspace*{-.5cm} }
%	\begin{framed}
%		\small
%	Wir verwenden die Physikkonvention, eine zeitliche Ableitung als Punkt darzustellen; $\dot{x} = \frac{dx\left(t\right)}{dt} = \frac{d}{dt}\left(x\left(t\right)\right)$.
%	\end{framed}
	

\thispagestyle{empty}


\noindent
\begin{Exercise}[label = ants, difficulty = 3, title = Ameisen im Dreieck]
	Drei Ameisen sitzten an den Eckpunkten eines gleichseiten Dreiecks. Jede Ameise bewegt sich direkt auf ihren jeweils rechten Nachbarn mit der Geschwindigkeit $v$ zu. Stelle eine Gleichung der Form $r = f\left(\varphi\right)$ auf, die den Abstand $r$ einer Ameise vom Treffpunkt der drei in Abhängigkeit eines \textit{sinnvoll} definierten Winkels $\varphi$ angibt.
\end{Exercise}
\begin{Exercise}[label = rocket1, title = Rakete im Weltraum, difficulty = 4, origin = Aaron Wild]
	Eine Rakete bewegt sich im Weltraum. Dabei soll sie am Anfang in Ruhe sein. Dann wird das Triebwerk gezündet, welches eine konstante Auströmgeschwindigkeit $v_a$ und mit konstantem Treibstoffverbrauch $|\dot{m}| = \lambda =\mathrm{const.}$ betrieben wird. \\
	Bei vollem Tank hat die Rakete eine Masse $m_0$ und bei leerem eine Masse $m_1<m_0$. Welche Geschwindigkeit hat die Rakete erreicht, wenn der Tank leer ist?
\end{Exercise} 
\begin{Exercise}[label = halfplan, title = Ein halber Planet, origin = {4.Runde IPhO 2015, Kurzaufgabe}, difficulty = 3]
	Auf einem halbkugelförmigen Planeten beträgt die Schwerebeschleunigung direkt in der Mitte der flachen Deckfläche $g_0$.  Die Dichte ist im gesamten Planeten gleich, und beträgt $\rho$. Wie groß ist der Planetenradius?
\end{Exercise}



\end{document}
